% Options for packages loaded elsewhere
% Options for packages loaded elsewhere
\PassOptionsToPackage{unicode}{hyperref}
\PassOptionsToPackage{hyphens}{url}
\PassOptionsToPackage{dvipsnames,svgnames,x11names}{xcolor}
%
\documentclass[
  letterpaper,
  DIV=11,
  numbers=noendperiod]{scrartcl}
\usepackage{xcolor}
\usepackage{amsmath,amssymb}
\setcounter{secnumdepth}{-\maxdimen} % remove section numbering
\usepackage{iftex}
\ifPDFTeX
  \usepackage[T1]{fontenc}
  \usepackage[utf8]{inputenc}
  \usepackage{textcomp} % provide euro and other symbols
\else % if luatex or xetex
  \usepackage{unicode-math} % this also loads fontspec
  \defaultfontfeatures{Scale=MatchLowercase}
  \defaultfontfeatures[\rmfamily]{Ligatures=TeX,Scale=1}
\fi
\usepackage{lmodern}
\ifPDFTeX\else
  % xetex/luatex font selection
\fi
% Use upquote if available, for straight quotes in verbatim environments
\IfFileExists{upquote.sty}{\usepackage{upquote}}{}
\IfFileExists{microtype.sty}{% use microtype if available
  \usepackage[]{microtype}
  \UseMicrotypeSet[protrusion]{basicmath} % disable protrusion for tt fonts
}{}
\makeatletter
\@ifundefined{KOMAClassName}{% if non-KOMA class
  \IfFileExists{parskip.sty}{%
    \usepackage{parskip}
  }{% else
    \setlength{\parindent}{0pt}
    \setlength{\parskip}{6pt plus 2pt minus 1pt}}
}{% if KOMA class
  \KOMAoptions{parskip=half}}
\makeatother
% Make \paragraph and \subparagraph free-standing
\makeatletter
\ifx\paragraph\undefined\else
  \let\oldparagraph\paragraph
  \renewcommand{\paragraph}{
    \@ifstar
      \xxxParagraphStar
      \xxxParagraphNoStar
  }
  \newcommand{\xxxParagraphStar}[1]{\oldparagraph*{#1}\mbox{}}
  \newcommand{\xxxParagraphNoStar}[1]{\oldparagraph{#1}\mbox{}}
\fi
\ifx\subparagraph\undefined\else
  \let\oldsubparagraph\subparagraph
  \renewcommand{\subparagraph}{
    \@ifstar
      \xxxSubParagraphStar
      \xxxSubParagraphNoStar
  }
  \newcommand{\xxxSubParagraphStar}[1]{\oldsubparagraph*{#1}\mbox{}}
  \newcommand{\xxxSubParagraphNoStar}[1]{\oldsubparagraph{#1}\mbox{}}
\fi
\makeatother

\usepackage{color}
\usepackage{fancyvrb}
\newcommand{\VerbBar}{|}
\newcommand{\VERB}{\Verb[commandchars=\\\{\}]}
\DefineVerbatimEnvironment{Highlighting}{Verbatim}{commandchars=\\\{\}}
% Add ',fontsize=\small' for more characters per line
\usepackage{framed}
\definecolor{shadecolor}{RGB}{241,243,245}
\newenvironment{Shaded}{\begin{snugshade}}{\end{snugshade}}
\newcommand{\AlertTok}[1]{\textcolor[rgb]{0.68,0.00,0.00}{#1}}
\newcommand{\AnnotationTok}[1]{\textcolor[rgb]{0.37,0.37,0.37}{#1}}
\newcommand{\AttributeTok}[1]{\textcolor[rgb]{0.40,0.45,0.13}{#1}}
\newcommand{\BaseNTok}[1]{\textcolor[rgb]{0.68,0.00,0.00}{#1}}
\newcommand{\BuiltInTok}[1]{\textcolor[rgb]{0.00,0.23,0.31}{#1}}
\newcommand{\CharTok}[1]{\textcolor[rgb]{0.13,0.47,0.30}{#1}}
\newcommand{\CommentTok}[1]{\textcolor[rgb]{0.37,0.37,0.37}{#1}}
\newcommand{\CommentVarTok}[1]{\textcolor[rgb]{0.37,0.37,0.37}{\textit{#1}}}
\newcommand{\ConstantTok}[1]{\textcolor[rgb]{0.56,0.35,0.01}{#1}}
\newcommand{\ControlFlowTok}[1]{\textcolor[rgb]{0.00,0.23,0.31}{\textbf{#1}}}
\newcommand{\DataTypeTok}[1]{\textcolor[rgb]{0.68,0.00,0.00}{#1}}
\newcommand{\DecValTok}[1]{\textcolor[rgb]{0.68,0.00,0.00}{#1}}
\newcommand{\DocumentationTok}[1]{\textcolor[rgb]{0.37,0.37,0.37}{\textit{#1}}}
\newcommand{\ErrorTok}[1]{\textcolor[rgb]{0.68,0.00,0.00}{#1}}
\newcommand{\ExtensionTok}[1]{\textcolor[rgb]{0.00,0.23,0.31}{#1}}
\newcommand{\FloatTok}[1]{\textcolor[rgb]{0.68,0.00,0.00}{#1}}
\newcommand{\FunctionTok}[1]{\textcolor[rgb]{0.28,0.35,0.67}{#1}}
\newcommand{\ImportTok}[1]{\textcolor[rgb]{0.00,0.46,0.62}{#1}}
\newcommand{\InformationTok}[1]{\textcolor[rgb]{0.37,0.37,0.37}{#1}}
\newcommand{\KeywordTok}[1]{\textcolor[rgb]{0.00,0.23,0.31}{\textbf{#1}}}
\newcommand{\NormalTok}[1]{\textcolor[rgb]{0.00,0.23,0.31}{#1}}
\newcommand{\OperatorTok}[1]{\textcolor[rgb]{0.37,0.37,0.37}{#1}}
\newcommand{\OtherTok}[1]{\textcolor[rgb]{0.00,0.23,0.31}{#1}}
\newcommand{\PreprocessorTok}[1]{\textcolor[rgb]{0.68,0.00,0.00}{#1}}
\newcommand{\RegionMarkerTok}[1]{\textcolor[rgb]{0.00,0.23,0.31}{#1}}
\newcommand{\SpecialCharTok}[1]{\textcolor[rgb]{0.37,0.37,0.37}{#1}}
\newcommand{\SpecialStringTok}[1]{\textcolor[rgb]{0.13,0.47,0.30}{#1}}
\newcommand{\StringTok}[1]{\textcolor[rgb]{0.13,0.47,0.30}{#1}}
\newcommand{\VariableTok}[1]{\textcolor[rgb]{0.07,0.07,0.07}{#1}}
\newcommand{\VerbatimStringTok}[1]{\textcolor[rgb]{0.13,0.47,0.30}{#1}}
\newcommand{\WarningTok}[1]{\textcolor[rgb]{0.37,0.37,0.37}{\textit{#1}}}

\usepackage{longtable,booktabs,array}
\newcounter{none} % for unnumbered tables
\usepackage{calc} % for calculating minipage widths
% Correct order of tables after \paragraph or \subparagraph
\usepackage{etoolbox}
\makeatletter
\patchcmd\longtable{\par}{\if@noskipsec\mbox{}\fi\par}{}{}
\makeatother
% Allow footnotes in longtable head/foot
\IfFileExists{footnotehyper.sty}{\usepackage{footnotehyper}}{\usepackage{footnote}}
\makesavenoteenv{longtable}
\usepackage{graphicx}
\makeatletter
\newsavebox\pandoc@box
\newcommand*\pandocbounded[1]{% scales image to fit in text height/width
  \sbox\pandoc@box{#1}%
  \Gscale@div\@tempa{\textheight}{\dimexpr\ht\pandoc@box+\dp\pandoc@box\relax}%
  \Gscale@div\@tempb{\linewidth}{\wd\pandoc@box}%
  \ifdim\@tempb\p@<\@tempa\p@\let\@tempa\@tempb\fi% select the smaller of both
  \ifdim\@tempa\p@<\p@\scalebox{\@tempa}{\usebox\pandoc@box}%
  \else\usebox{\pandoc@box}%
  \fi%
}
% Set default figure placement to htbp
\def\fps@figure{htbp}
\makeatother





\setlength{\emergencystretch}{3em} % prevent overfull lines

\providecommand{\tightlist}{%
  \setlength{\itemsep}{0pt}\setlength{\parskip}{0pt}}



 


\KOMAoption{captions}{tableheading}
\makeatletter
\@ifpackageloaded{caption}{}{\usepackage{caption}}
\AtBeginDocument{%
\ifdefined\contentsname
  \renewcommand*\contentsname{Table of contents}
\else
  \newcommand\contentsname{Table of contents}
\fi
\ifdefined\listfigurename
  \renewcommand*\listfigurename{List of Figures}
\else
  \newcommand\listfigurename{List of Figures}
\fi
\ifdefined\listtablename
  \renewcommand*\listtablename{List of Tables}
\else
  \newcommand\listtablename{List of Tables}
\fi
\ifdefined\figurename
  \renewcommand*\figurename{Figure}
\else
  \newcommand\figurename{Figure}
\fi
\ifdefined\tablename
  \renewcommand*\tablename{Table}
\else
  \newcommand\tablename{Table}
\fi
}
\@ifpackageloaded{float}{}{\usepackage{float}}
\floatstyle{ruled}
\@ifundefined{c@chapter}{\newfloat{codelisting}{h}{lop}}{\newfloat{codelisting}{h}{lop}[chapter]}
\floatname{codelisting}{Listing}
\newcommand*\listoflistings{\listof{codelisting}{List of Listings}}
\makeatother
\makeatletter
\makeatother
\makeatletter
\@ifpackageloaded{caption}{}{\usepackage{caption}}
\@ifpackageloaded{subcaption}{}{\usepackage{subcaption}}
\makeatother
\usepackage{bookmark}
\IfFileExists{xurl.sty}{\usepackage{xurl}}{} % add URL line breaks if available
\urlstyle{same}
% fallback for those not using the hyperref driver hyperxmp:
\makeatletter
\@ifundefined{xmpquote}{\newcommand{\xmpquote}[1]{#1}}{}
\makeatother
\hypersetup{
  pdftitle={Steam Games Analysis},
  pdfauthor={Howard Song, Milan Bertolutti, Yuanchen Ling, Maryamalsadat Hosseini},
  colorlinks=true,
  linkcolor={blue},
  filecolor={Maroon},
  citecolor={Blue},
  urlcolor={Blue},
  pdfcreator={LaTeX via pandoc}}


\title{Steam Games Analysis}
\author{Howard Song, Milan Bertolutti, Yuanchen Ling, Maryamalsadat
Hosseini}
\date{}
\begin{document}
\maketitle


\section{Summary}\label{summary}

This project investigates how well game characteristics can predict the
total number of reviews received by Steam games released between 2018
and 2022. Using the Steam Games Dataset, we explore relationships
between review counts and game features, including price, genres,
categories, number of supported languages, and release date. We conduct
exploratory data analysis to examine review distributions and feature
relationships, followed by regression modeling to evaluate predictive
performance. The results show that our linear regression model explains
approximately one-third of the variation in log review counts on the
test set and performs better than the mean baseline in cross-validation.
These findings suggest that game characteristics provide useful
predictive information and may help developers and publishers better
understand patterns in review activity. However, the model is less
accurate for games with very few or very high review counts, and other
factors, such as marketing activities, developer reputation, and
community engagement, are not included in the model.

\section{Introduction}\label{introduction}

The video game industry has become a major part of the digital
entertainment market, with online platforms such as Steam providing
developers and publishers access to a global audience. In addition to
purchasing and playing games, Steam users can share their experiences
and recommend games through reviews (Valve, n.d.). These reviews provide
insights into player opinions and serve as a measurable indicator of
audience engagement.

Understanding the factors associated with review volume is important for
game developers and publishers because reviews can provide insights into
player engagement and a game's visibility among potential players.
Previous research has examined the relationship between online reviews
and video game usage (Philp \& Nepomuceno, 2024), as well as how
different player engagement modes influence review characteristics
(Guzsvinecz \& Szűcs, 2026). Research on the Steam platform has also
investigated game reviews and their relationship with game
characteristics (Lin et al., 2019). These studies provide relevant
background for understanding review activity in video games. However,
examining how selected game characteristics can be used to predict the
total number of reviews remains the focus of our study. Characteristics
such as price, genres, available features, language support, and release
timing may be associated with differences in the number of reviews
received. Examining these relationships may help developers and
publishers better understand which game characteristics are associated
with greater audience engagement.

This study investigates the following research question: How well can
selected game characteristics predict the total number of reviews
received by Steam games released between 2018 and 2022? Specifically,
the analysis examines five groups of explanatory variables: game price,
number of supported languages, release date, genres, and Steam
categories. Total review count, calculated by adding positive and
negative reviews, is used as the response variable.

The analysis uses the Steam Games Dataset published by FronkonGames on
Hugging Face (FronkonGames, n.d.). Each row represents an individual
Steam game, with information including release date, price, genres,
categories, language support, and positive and negative review counts.
The study focuses on 46,859 games released between 2018 and 2022. The
dataset is randomly divided into training and test sets, with 80\% of
the data used for exploratory data analysis and model development and
the remaining 20\% reserved for final model evaluation.

To investigate the research question, we first conduct exploratory data
analysis to examine the distribution of review counts and explore their
associations with selected game characteristics. We then develop a
linear regression model to predict the logarithm of total review counts
and compare its performance with a baseline model that predicts the
mean. Model performance is evaluated using five-fold cross-validation on
the training set and a held-out test set.

This study aims to evaluate the predictive value of selected game
characteristics rather than establish causal relationships. By examining
both feature associations and predictive performance, the analysis
provides insights into how well game metadata can explain and predict
differences in review volume, while highlighting the limitations of
using these characteristics alone to estimate player engagement.

\section{Methods}\label{methods}

\subsection{Data}\label{data}

The code below downloads the Steam Games Dataset (FronkonGames, n.d.)
from Hugging Face and saves an unmodified copy to \texttt{data/raw/}.
The download is skipped if that file already exists. We then keep games
released between 2018 and 2022, drop rows with missing review counts,
duplicate games or unreadable release dates, and save the cleaned data
to \texttt{data/processed/}. This leaves 46,859 games, which we split at
random into a training set (80\%) and a test set (20\%) using seed 1234.
The split happens before any exploratory analysis: only the training set
is used for exploration and model building, and the test set is used
once, at the end, to report final performance.

\begin{Shaded}
\begin{Highlighting}[]
\ImportTok{import}\NormalTok{ sys}
\ImportTok{from}\NormalTok{ pathlib }\ImportTok{import}\NormalTok{ Path}

\ImportTok{import}\NormalTok{ pandas }\ImportTok{as}\NormalTok{ pd}
\ImportTok{from}\NormalTok{ sklearn.model\_selection }\ImportTok{import}\NormalTok{ train\_test\_split}

\CommentTok{\# Reuse the download and cleaning functions from scripts/process\_data.py}
\NormalTok{sys.path.append(}\StringTok{"../scripts"}\NormalTok{)}
\ImportTok{from}\NormalTok{ process\_data }\ImportTok{import}\NormalTok{ (}
\NormalTok{    DATA\_URL,}
\NormalTok{    EXPECTED\_COHORT\_SIZE,}
\NormalTok{    RAW\_FILENAME,}
\NormalTok{    SEED,}
\NormalTok{    TEST\_SIZE,}
\NormalTok{    clean\_data,}
\NormalTok{    download\_data,}
\NormalTok{)}

\NormalTok{raw\_dir }\OperatorTok{=}\NormalTok{ Path(}\StringTok{"../data/raw"}\NormalTok{)}
\NormalTok{processed\_dir }\OperatorTok{=}\NormalTok{ Path(}\StringTok{"../data/processed"}\NormalTok{)}
\NormalTok{processed\_dir.mkdir(parents}\OperatorTok{=}\VariableTok{True}\NormalTok{, exist\_ok}\OperatorTok{=}\VariableTok{True}\NormalTok{)}

\CommentTok{\# 1. Download the raw data (skipped if it is already in data/raw/)}
\NormalTok{raw\_path }\OperatorTok{=}\NormalTok{ download\_data(DATA\_URL, raw\_dir }\OperatorTok{/}\NormalTok{ RAW\_FILENAME)}
\NormalTok{raw }\OperatorTok{=}\NormalTok{ pd.read\_parquet(raw\_path)}

\CommentTok{\# 2. Clean it and record every dropped row with the reason}
\NormalTok{cleaned, ledger }\OperatorTok{=}\NormalTok{ clean\_data(raw)}
\ControlFlowTok{if} \BuiltInTok{len}\NormalTok{(cleaned) }\OperatorTok{!=}\NormalTok{ EXPECTED\_COHORT\_SIZE:}
    \ControlFlowTok{raise} \PreprocessorTok{ValueError}\NormalTok{(}
        \SpecialStringTok{f"Expected }\SpecialCharTok{\{}\NormalTok{EXPECTED\_COHORT\_SIZE}\SpecialCharTok{:,\}}\SpecialStringTok{ games after cleaning, got }\SpecialCharTok{\{}\BuiltInTok{len}\NormalTok{(cleaned)}\SpecialCharTok{:,\}}\SpecialStringTok{."}
\NormalTok{    )}

\CommentTok{\# 3. Split before any exploration}
\NormalTok{train\_split, test\_split }\OperatorTok{=}\NormalTok{ train\_test\_split(}
\NormalTok{    cleaned, test\_size}\OperatorTok{=}\NormalTok{TEST\_SIZE, random\_state}\OperatorTok{=}\NormalTok{SEED}
\NormalTok{)}

\NormalTok{cleaned.to\_csv(processed\_dir }\OperatorTok{/} \StringTok{"cleaned.csv"}\NormalTok{, index}\OperatorTok{=}\VariableTok{False}\NormalTok{)}
\NormalTok{ledger.to\_csv(processed\_dir }\OperatorTok{/} \StringTok{"drop{-}ledger.csv"}\NormalTok{, index}\OperatorTok{=}\VariableTok{False}\NormalTok{)}
\NormalTok{train\_split.to\_csv(processed\_dir }\OperatorTok{/} \StringTok{"train.csv"}\NormalTok{, index}\OperatorTok{=}\VariableTok{False}\NormalTok{)}
\NormalTok{test\_split.to\_csv(processed\_dir }\OperatorTok{/} \StringTok{"test.csv"}\NormalTok{, index}\OperatorTok{=}\VariableTok{False}\NormalTok{)}

\NormalTok{ledger[}\StringTok{"reason"}\NormalTok{].value\_counts()}
\end{Highlighting}
\end{Shaded}

\begin{verbatim}
Downloading https://huggingface.co/datasets/FronkonGames/steam-games-dataset/resolve/6c21457b9ebf0a6a7f5b91e11281d65d456ea30e/data/train-00000-of-00001.parquet
  -> ../data/raw/steam_games.parquet
\end{verbatim}

\begin{verbatim}
reason
released outside 2018-2022    77287
Name: count, dtype: int64
\end{verbatim}

\begin{Shaded}
\begin{Highlighting}[]
\NormalTok{train }\OperatorTok{=}\NormalTok{ pd.read\_csv(processed\_dir }\OperatorTok{/} \StringTok{"train.csv"}\NormalTok{)}
\NormalTok{train.shape}
\end{Highlighting}
\end{Shaded}

\begin{verbatim}
(37487, 9)
\end{verbatim}

\subsection{Target variable}\label{target-variable}

We predict \texttt{log1p(total\_reviews)}, the log of reviews + 1.
Logarithmically transforming prevents the few very popular games from
dominating the model. Adding 1 to the value keeps games with zero
reviews defined.

\begin{Shaded}
\begin{Highlighting}[]
\ImportTok{import}\NormalTok{ numpy }\ImportTok{as}\NormalTok{ np}

\NormalTok{X\_train }\OperatorTok{=}\NormalTok{ train.drop(columns}\OperatorTok{=}\NormalTok{[}\StringTok{"total\_reviews"}\NormalTok{])}
\NormalTok{y\_train }\OperatorTok{=}\NormalTok{ np.log1p(train[}\StringTok{"total\_reviews"}\NormalTok{])}
\end{Highlighting}
\end{Shaded}

\subsection{Features}\label{features}

We used five groups of features, all built from the training data:

\begin{itemize}
\tightlist
\item
  \textbf{Price} in USD, used as is.
\item
  \textbf{Number of supported languages}, capped at 10, since adding
  languages beyond that point makes little further difference.
\item
  \textbf{Release day}: days since 1 January 2018, so the model can pick
  up trends over the 2018--2022 period.
\item
  \textbf{Genres} and \textbf{categories}: one yes/no column per genre
  and per Steam category. Only levels that appear in at least 1\% of
  training games are kept, so rare labels do not add noise.
  \texttt{Steam\ Trading\ Cards} is excluded because Valve grants it to
  games that already sell well, making it a consequence of popularity
  rather than a cause.
\end{itemize}

\texttt{appID} and \texttt{name} are identifiers and are not used as
features.

\subsubsection{Transformation}\label{transformation}

\begin{Shaded}
\begin{Highlighting}[]
\CommentTok{\# Transform the data to fit contract. Clip number of languages to a max of 10, and calculate release day.}
\NormalTok{X\_train[}\StringTok{"n\_languages"}\NormalTok{] }\OperatorTok{=}\NormalTok{ X\_train[}\StringTok{"n\_languages"}\NormalTok{].clip(upper}\OperatorTok{=}\DecValTok{10}\NormalTok{)}
\NormalTok{X\_train[}\StringTok{"release\_day"}\NormalTok{] }\OperatorTok{=}\NormalTok{ (pd.to\_datetime(X\_train[}\StringTok{"release\_date"}\NormalTok{]) }\OperatorTok{{-}}\NormalTok{ pd.Timestamp(}\StringTok{"2018{-}01{-}01"}\NormalTok{)).dt.days}
\end{Highlighting}
\end{Shaded}

\subsubsection{Preprocessing}\label{preprocessing}

\begin{Shaded}
\begin{Highlighting}[]
\ImportTok{from}\NormalTok{ sklearn.compose }\ImportTok{import}\NormalTok{ ColumnTransformer}
\ImportTok{from}\NormalTok{ sklearn.feature\_extraction.text }\ImportTok{import}\NormalTok{ CountVectorizer}
\ImportTok{from}\NormalTok{ sklearn.preprocessing }\ImportTok{import}\NormalTok{ StandardScaler}

\NormalTok{list\_encoder }\OperatorTok{=}\NormalTok{ CountVectorizer(}
\NormalTok{    token\_pattern}\OperatorTok{=}\VerbatimStringTok{r\textquotesingle{}"}\KeywordTok{(}\PreprocessorTok{[\^{}"]}\OperatorTok{+}\KeywordTok{)}\VerbatimStringTok{"\textquotesingle{}}\NormalTok{, }\CommentTok{\# This keeps multi word labels}
\NormalTok{    lowercase}\OperatorTok{=}\VariableTok{False}\NormalTok{,}
\NormalTok{    binary}\OperatorTok{=}\VariableTok{True}\NormalTok{,}
\NormalTok{    min\_df}\OperatorTok{=}\FloatTok{0.01}\NormalTok{,}
\NormalTok{    stop\_words}\OperatorTok{=}\NormalTok{[}\StringTok{"Steam Trading Cards"}\NormalTok{],}
\NormalTok{)}

\NormalTok{preprocessor }\OperatorTok{=}\NormalTok{ ColumnTransformer([}
\NormalTok{    (}\StringTok{"numeric"}\NormalTok{, StandardScaler(), [}\StringTok{"price"}\NormalTok{, }\StringTok{"n\_languages"}\NormalTok{, }\StringTok{"release\_day"}\NormalTok{]),}
\NormalTok{    (}\StringTok{"genre"}\NormalTok{, list\_encoder, }\StringTok{"genres"}\NormalTok{),}
\NormalTok{    (}\StringTok{"category"}\NormalTok{, list\_encoder, }\StringTok{"categories"}\NormalTok{),}
\NormalTok{])}
\end{Highlighting}
\end{Shaded}

\subsection{Model and evaluation}\label{model-and-evaluation}

We fit a linear regression, with the preprocessing and model combined in
one pipeline. We compare it to a baseline that always predicts the mean.

Both models are scored with 5-fold cross-validation on the training set.
Because preprocessing is inside the pipeline, it is re-fit on each fold,
so the validation data never influences it.

Our main metric is Root Mean Square Error (RMSE) on the \texttt{log1p}
scale, which measures relative error so that a few very popular games do
not dominate. We also report Mean Absolute Error (MAE) and R².

\begin{Shaded}
\begin{Highlighting}[]
\ImportTok{from}\NormalTok{ sklearn.dummy }\ImportTok{import}\NormalTok{ DummyRegressor}
\ImportTok{from}\NormalTok{ sklearn.linear\_model }\ImportTok{import}\NormalTok{ LinearRegression}
\ImportTok{from}\NormalTok{ sklearn.model\_selection }\ImportTok{import}\NormalTok{ cross\_validate}
\ImportTok{from}\NormalTok{ sklearn.pipeline }\ImportTok{import}\NormalTok{ make\_pipeline}

\NormalTok{pipe }\OperatorTok{=}\NormalTok{ make\_pipeline(preprocessor, LinearRegression())}

\NormalTok{models }\OperatorTok{=}\NormalTok{ \{}
    \StringTok{"Baseline (mean)"}\NormalTok{: DummyRegressor(),}
    \StringTok{"Linear regression"}\NormalTok{: pipe,}
\NormalTok{\}}

\NormalTok{scoring }\OperatorTok{=}\NormalTok{ \{}
    \StringTok{"RMSE"}\NormalTok{: }\StringTok{"neg\_root\_mean\_squared\_error"}\NormalTok{,}
    \StringTok{"MAE"}\NormalTok{: }\StringTok{"neg\_mean\_absolute\_error"}\NormalTok{,}
    \StringTok{"R²"}\NormalTok{: }\StringTok{"r2"}\NormalTok{,}
\NormalTok{\}}

\NormalTok{cv\_results }\OperatorTok{=}\NormalTok{ \{\}}
\ControlFlowTok{for}\NormalTok{ name, model }\KeywordTok{in}\NormalTok{ models.items():}
\NormalTok{    scores }\OperatorTok{=}\NormalTok{ cross\_validate(model, X\_train, y\_train, cv}\OperatorTok{=}\DecValTok{5}\NormalTok{, scoring}\OperatorTok{=}\NormalTok{scoring)}
\NormalTok{    cv\_results[name] }\OperatorTok{=}\NormalTok{ \{}
        \StringTok{"RMSE"}\NormalTok{: }\OperatorTok{{-}}\NormalTok{scores[}\StringTok{"test\_RMSE"}\NormalTok{].mean(),}
        \StringTok{"MAE"}\NormalTok{: }\OperatorTok{{-}}\NormalTok{scores[}\StringTok{"test\_MAE"}\NormalTok{].mean(),}
        \StringTok{"R²"}\NormalTok{: scores[}\StringTok{"test\_R²"}\NormalTok{].mean(),}
\NormalTok{    \}}

\NormalTok{cv\_results }\OperatorTok{=}\NormalTok{ pd.DataFrame(cv\_results).T.}\BuiltInTok{round}\NormalTok{(}\DecValTok{3}\NormalTok{)}
\end{Highlighting}
\end{Shaded}

\section{Results}\label{results}

\subsection{Exploratory Data Analysis}\label{exploratory-data-analysis}

\subsubsection{Summary Statistics}\label{summary-statistics}

The training dataset contains 37,487 Steam games. Descriptive statistics
for the numerical variables are presented in
Table~\ref{tbl-training-summary}.

\begin{longtable}[]{@{}lllllllll@{}}

\caption{\label{tbl-training-summary}Descriptive statistics of the
training dataset}

\tabularnewline

\toprule\noalign{}
& N & Mean & SD & Min & Q1 & Median & Q3 & Max \\
Variable & & & & & & & & \\
\midrule\noalign{}
\endhead
\bottomrule\noalign{}
\endlastfoot
Price (USD) & 37487 & 4.58 & 12.60 & 0.0 & 0.59 & 1.99 & 4.99 & 999.0 \\
Supported languages (count) & 37487 & 4.54 & 11.68 & 0.0 & 1.00 & 1.00 &
3.00 & 103.0 \\
Total reviews (count) & 37487 & 1158.54 & 17109.06 & 0.0 & 2.00 & 14.00
& 61.00 & 1227784.0 \\

\end{longtable}

The mean number of reviews (1,158.54) is substantially higher than the
median (14), indicating a strongly right-skewed distribution. This
motivates the log transformation used in the exploratory visualization.

The exploratory data analysis was conducted using only the training
dataset to avoid data leakage.

\subsubsection{Distribution of Total
Reviews}\label{distribution-of-total-reviews}

\begin{figure}

\centering{

\pandocbounded{\includegraphics[keepaspectratio]{videogames-regression_files/figure-pdf/fig-review-distribution-output-1.pdf}}

}

\caption{\label{fig-review-distribution}Distribution of log-transformed
total reviews among Steam games in the training dataset. The log(1 +
total reviews) transformation accounts for games with zero reviews and
reduces the influence of extreme values.}

\end{figure}%

The distribution of total reviews is highly right-skewed
(Figure~\ref{fig-review-distribution}). Most games receive relatively
few reviews, while a small number of games receive substantially more.
Applying the log transformation makes the distribution easier to
visualize while retaining games with zero reviews.

\subsubsection{Supported Languages and Review
Counts}\label{supported-languages-and-review-counts}

\begin{figure}

\centering{

\pandocbounded{\includegraphics[keepaspectratio]{videogames-regression_files/figure-pdf/fig-language-reviews-output-1.pdf}}

}

\caption{\label{fig-language-reviews}Median total reviews by the number
of supported languages. The sample size of each group is displayed above
the corresponding bar.}

\end{figure}%

Games supporting more languages tend to have higher median review counts
(Figure~\ref{fig-language-reviews}). This suggests a positive
association between language availability and review volume. However,
this relationship is descriptive and does not establish causality.

\subsection{Model performance}\label{model-performance}

After cross-validation, we fit the linear regression on the full
training set and scored it once on the test set.

\begin{Shaded}
\begin{Highlighting}[]
\ImportTok{from}\NormalTok{ sklearn.metrics }\ImportTok{import}\NormalTok{ root\_mean\_squared\_error, mean\_absolute\_error, r2\_score}

\NormalTok{pipe.fit(X\_train, y\_train)}

\CommentTok{\# Read and transform test data in same way as training data was}
\NormalTok{test }\OperatorTok{=}\NormalTok{ pd.read\_csv(}\StringTok{"../data/processed/test.csv"}\NormalTok{)}
\NormalTok{X\_test }\OperatorTok{=}\NormalTok{ test.drop(columns}\OperatorTok{=}\NormalTok{[}\StringTok{"total\_reviews"}\NormalTok{])}
\NormalTok{y\_test }\OperatorTok{=}\NormalTok{ np.log1p(test[}\StringTok{"total\_reviews"}\NormalTok{])}
\NormalTok{X\_test[}\StringTok{"n\_languages"}\NormalTok{] }\OperatorTok{=}\NormalTok{ X\_test[}\StringTok{"n\_languages"}\NormalTok{].clip(upper}\OperatorTok{=}\DecValTok{10}\NormalTok{)}
\NormalTok{X\_test[}\StringTok{"release\_day"}\NormalTok{] }\OperatorTok{=}\NormalTok{ (pd.to\_datetime(X\_test[}\StringTok{"release\_date"}\NormalTok{]) }\OperatorTok{{-}}\NormalTok{ pd.Timestamp(}\StringTok{"2018{-}01{-}01"}\NormalTok{)).dt.days}

\NormalTok{y\_pred }\OperatorTok{=}\NormalTok{ pipe.predict(X\_test)}

\NormalTok{test\_results }\OperatorTok{=}\NormalTok{ pd.DataFrame(\{}
    \StringTok{"RMSE"}\NormalTok{: [root\_mean\_squared\_error(y\_test, y\_pred)],}
    \StringTok{"MAE"}\NormalTok{: [mean\_absolute\_error(y\_test, y\_pred)],}
    \StringTok{"R²"}\NormalTok{: [r2\_score(y\_test, y\_pred)],}
\NormalTok{\}, index}\OperatorTok{=}\NormalTok{[}\StringTok{"Linear regression"}\NormalTok{]).}\BuiltInTok{round}\NormalTok{(}\DecValTok{3}\NormalTok{)}

\CommentTok{\# Numbers quoted in the text below}
\NormalTok{cv\_base }\OperatorTok{=}\NormalTok{ cv\_results.loc[}\StringTok{"Baseline (mean)"}\NormalTok{]}
\NormalTok{cv\_lr }\OperatorTok{=}\NormalTok{ cv\_results.loc[}\StringTok{"Linear regression"}\NormalTok{]}
\NormalTok{test\_lr }\OperatorTok{=}\NormalTok{ test\_results.loc[}\StringTok{"Linear regression"}\NormalTok{]}
\NormalTok{error\_factor }\OperatorTok{=}\NormalTok{ np.exp(test\_lr[}\StringTok{"MAE"}\NormalTok{])  }\CommentTok{\# MAE on the log scale {-}\textgreater{} multiplicative error}
\end{Highlighting}
\end{Shaded}

\begin{Shaded}
\begin{Highlighting}[]
\NormalTok{pd.concat(}
\NormalTok{    \{}\StringTok{"Cross{-}validation"}\NormalTok{: cv\_results, }\StringTok{"Test set"}\NormalTok{: test\_results\},}
\NormalTok{    names}\OperatorTok{=}\NormalTok{[}\StringTok{"Evaluated on"}\NormalTok{, }\StringTok{"Model"}\NormalTok{],}
\NormalTok{) }\OperatorTok{+} \FloatTok{0.0}  \CommentTok{\# turns the baseline\textquotesingle{}s rounded R² of {-}0.000 into 0.000}
\end{Highlighting}
\end{Shaded}

\begin{longtable}[]{@{}lllll@{}}

\caption{\label{tbl-performance}}

\tabularnewline

\toprule\noalign{}
& & RMSE & MAE & R² \\
Evaluated on & Model & & & \\
\midrule\noalign{}
\endhead
\bottomrule\noalign{}
\endlastfoot
\multirow{2}{*}{Cross-validation} & Baseline (mean) & 2.289 & 1.778 &
0.000 \\
& Linear regression & 1.846 & 1.436 & 0.349 \\
Test set & Linear regression & 1.829 & 1.417 & 0.340 \\

\end{longtable}

Linear regression beats the mean baseline on every metric
(Table~\ref{tbl-performance}). In cross-validation it lowers RMSE from
2.29 to 1.85 and explains 35\% of the variance in log reviews, while the
baseline explains none. The improvements are their but the model only
explains some of the variance. The features Price, number of languages,
release day, genres and categories together account for about a third of
the variation in how many reviews a game gets. This indicates that there
are features out there that explain more of that variance that was not
trained in the model.

The gap between the training validation is very close together (RMSE
1.83 vs 1.85; R² 0.34 vs 0.35). This means the model is not overfitting
and should perform about as well on other Steam games released in 2018 -
2022. In practical terms, the test MAE is 1.42 on the log scale. Because
the model predicts \(\log(1 + y)\), a gap of \(d\) between the actual
count \(y\) and the predicted count \(\hat{y}\) on the log scale is a
ratio of \(e^{d}\) on the original scale:

\[
\log(1 + y) - \log(1 + \hat{y}) = d
\quad \Longleftrightarrow \quad
\frac{1 + y}{1 + \hat{y}} = e^{d}
\]

Setting \(d\) to the MAE gives \(e^{\text{MAE}} \approx\) 4.1, so a
typical prediction is off by a factor of about 4.1 in either direction.
For example, a game predicted to receive 100 reviews would typically end
up with somewhere between 24 and 412.

\subsection{Prediction accuracy by game
popularity}\label{prediction-accuracy-by-game-popularity}

\subsubsection{Visualization Code:}\label{visualization-code}

\begin{Shaded}
\begin{Highlighting}[]
\ImportTok{import}\NormalTok{ altair }\ImportTok{as}\NormalTok{ alt}

\CommentTok{\# \textgreater{} Citation Claude Opus 5.5 was used for this visualization}

\CommentTok{\# Group test games by their actual number of reviews.}
\CommentTok{\# Groups are about half a power of ten wide, so they are evenly spaced on a log scale.}
\NormalTok{edges }\OperatorTok{=}\NormalTok{ [}\OperatorTok{{-}}\DecValTok{1}\NormalTok{, }\DecValTok{0}\NormalTok{, }\DecValTok{3}\NormalTok{, }\DecValTok{10}\NormalTok{, }\DecValTok{30}\NormalTok{, }\DecValTok{100}\NormalTok{, }\DecValTok{300}\NormalTok{, }\DecValTok{1000}\NormalTok{, }\DecValTok{3000}\NormalTok{, np.inf]}
\NormalTok{review\_group }\OperatorTok{=}\NormalTok{ pd.cut(}
\NormalTok{    test[}\StringTok{"total\_reviews"}\NormalTok{],}
\NormalTok{    bins}\OperatorTok{=}\NormalTok{edges,}
\NormalTok{    labels}\OperatorTok{=}\NormalTok{[}\StringTok{"0"}\NormalTok{, }\StringTok{"1–3"}\NormalTok{, }\StringTok{"4–10"}\NormalTok{, }\StringTok{"11–30"}\NormalTok{, }\StringTok{"31–100"}\NormalTok{, }\StringTok{"101–300"}\NormalTok{, }\StringTok{"301–1,000"}\NormalTok{, }\StringTok{"1,001–3,000"}\NormalTok{, }\StringTok{"3,001+"}\NormalTok{],}
\NormalTok{)}

\NormalTok{predictions }\OperatorTok{=}\NormalTok{ pd.DataFrame(\{}
    \StringTok{"group"}\NormalTok{: review\_group,}
    \StringTok{"actual"}\NormalTok{: test[}\StringTok{"total\_reviews"}\NormalTok{],}
    \StringTok{"predicted"}\NormalTok{: np.expm1(y\_pred),  }\CommentTok{\# back to review counts}
\NormalTok{\})}

\NormalTok{group\_summary }\OperatorTok{=}\NormalTok{ predictions.groupby(}\StringTok{"group"}\NormalTok{, observed}\OperatorTok{=}\VariableTok{True}\NormalTok{).agg(}
\NormalTok{    games}\OperatorTok{=}\NormalTok{(}\StringTok{"actual"}\NormalTok{, }\StringTok{"size"}\NormalTok{),}
\NormalTok{    actual}\OperatorTok{=}\NormalTok{(}\StringTok{"actual"}\NormalTok{, }\StringTok{"median"}\NormalTok{),}
\NormalTok{    predicted}\OperatorTok{=}\NormalTok{(}\StringTok{"predicted"}\NormalTok{, }\StringTok{"median"}\NormalTok{),}
\NormalTok{    predicted\_low}\OperatorTok{=}\NormalTok{(}\StringTok{"predicted"}\NormalTok{, }\KeywordTok{lambda}\NormalTok{ p: p.quantile(}\FloatTok{0.25}\NormalTok{)),}
\NormalTok{    predicted\_high}\OperatorTok{=}\NormalTok{(}\StringTok{"predicted"}\NormalTok{, }\KeywordTok{lambda}\NormalTok{ p: p.quantile(}\FloatTok{0.75}\NormalTok{)),}
\NormalTok{)}

\CommentTok{\# Numbers quoted in the text below}
\NormalTok{lowest }\OperatorTok{=}\NormalTok{ group\_summary.iloc[}\DecValTok{0}\NormalTok{]}
\NormalTok{highest }\OperatorTok{=}\NormalTok{ group\_summary.iloc[}\OperatorTok{{-}}\DecValTok{1}\NormalTok{]}

\CommentTok{\#| code{-}fold: true}
\NormalTok{alt.data\_transformers.disable\_max\_rows()  }\CommentTok{\# plot every test game (Altair\textquotesingle{}s default limit is 5,000 rows)}

\CommentTok{\# One shared scale for every review axis, so the two panels line up.}
\CommentTok{\# symlog\textquotesingle{}s default ticks bunch at the top, so place them at powers of ten.}
\NormalTok{max\_reviews }\OperatorTok{=}\NormalTok{ predictions[[}\StringTok{"actual"}\NormalTok{, }\StringTok{"predicted"}\NormalTok{]].}\BuiltInTok{max}\NormalTok{().}\BuiltInTok{max}\NormalTok{()}
\NormalTok{scale }\OperatorTok{=}\NormalTok{ alt.Scale(}\BuiltInTok{type}\OperatorTok{=}\StringTok{"symlog"}\NormalTok{, domain}\OperatorTok{=}\NormalTok{[}\DecValTok{0}\NormalTok{, max\_reviews])}
\NormalTok{axis }\OperatorTok{=}\NormalTok{ alt.Axis(values}\OperatorTok{=}\NormalTok{[}\DecValTok{0}\NormalTok{, }\DecValTok{1}\NormalTok{, }\DecValTok{10}\NormalTok{, }\DecValTok{100}\NormalTok{, }\DecValTok{1\_000}\NormalTok{, }\DecValTok{10\_000}\NormalTok{, }\DecValTok{100\_000}\NormalTok{, }\DecValTok{1\_000\_000}\NormalTok{],}
\NormalTok{                labelExpr}\OperatorTok{=}\StringTok{"format(datum.value, \textquotesingle{},\textquotesingle{})"}\NormalTok{)}

\CommentTok{\# (A) Every game, with the perfect{-}prediction line and the group boundaries}
\NormalTok{scatter }\OperatorTok{=}\NormalTok{ alt.Chart(predictions[[}\StringTok{"actual"}\NormalTok{, }\StringTok{"predicted"}\NormalTok{]]).mark\_circle(}
\NormalTok{    size}\OperatorTok{=}\DecValTok{8}\NormalTok{, opacity}\OperatorTok{=}\FloatTok{0.15}\NormalTok{, color}\OperatorTok{=}\StringTok{"\#1f77b4"}
\NormalTok{).encode(}
\NormalTok{    x}\OperatorTok{=}\NormalTok{alt.X(}\StringTok{"actual"}\NormalTok{, scale}\OperatorTok{=}\NormalTok{scale, axis}\OperatorTok{=}\NormalTok{axis, title}\OperatorTok{=}\StringTok{"Actual reviews"}\NormalTok{),}
\NormalTok{    y}\OperatorTok{=}\NormalTok{alt.Y(}\StringTok{"predicted"}\NormalTok{, scale}\OperatorTok{=}\NormalTok{scale, axis}\OperatorTok{=}\NormalTok{axis, title}\OperatorTok{=}\StringTok{"Predicted reviews"}\NormalTok{),}
\NormalTok{)}
\NormalTok{perfect }\OperatorTok{=}\NormalTok{ alt.Chart(pd.DataFrame(\{}\StringTok{"reviews"}\NormalTok{: [}\DecValTok{0}\NormalTok{, max\_reviews]\})).mark\_line(}
\NormalTok{    color}\OperatorTok{=}\StringTok{"black"}\NormalTok{, strokeDash}\OperatorTok{=}\NormalTok{[}\DecValTok{4}\NormalTok{, }\DecValTok{4}\NormalTok{]}
\NormalTok{).encode(}
\NormalTok{    x}\OperatorTok{=}\NormalTok{alt.X(}\StringTok{"reviews"}\NormalTok{, scale}\OperatorTok{=}\NormalTok{scale),}
\NormalTok{    y}\OperatorTok{=}\NormalTok{alt.Y(}\StringTok{"reviews"}\NormalTok{, scale}\OperatorTok{=}\NormalTok{scale),}
\NormalTok{)}
\NormalTok{boundaries }\OperatorTok{=}\NormalTok{ alt.Chart(pd.DataFrame(\{}\StringTok{"edge"}\NormalTok{: [e }\OperatorTok{+} \FloatTok{0.5} \ControlFlowTok{for}\NormalTok{ e }\KeywordTok{in}\NormalTok{ edges[}\DecValTok{1}\NormalTok{:}\OperatorTok{{-}}\DecValTok{1}\NormalTok{]]\})).mark\_rule(}
\NormalTok{    color}\OperatorTok{=}\StringTok{"grey"}\NormalTok{, opacity}\OperatorTok{=}\FloatTok{0.4}
\NormalTok{).encode(x}\OperatorTok{=}\NormalTok{alt.X(}\StringTok{"edge"}\NormalTok{, scale}\OperatorTok{=}\NormalTok{scale))}

\NormalTok{panel\_a }\OperatorTok{=}\NormalTok{ (boundaries }\OperatorTok{+}\NormalTok{ scatter }\OperatorTok{+}\NormalTok{ perfect).properties(}
\NormalTok{    width}\OperatorTok{=}\DecValTok{300}\NormalTok{, height}\OperatorTok{=}\DecValTok{300}\NormalTok{, title}\OperatorTok{=}\StringTok{"A. Every test game"}
\NormalTok{)}

\CommentTok{\# (B) Group medians, with the middle 50\% of predictions}
\NormalTok{summary }\OperatorTok{=}\NormalTok{ group\_summary.reset\_index()}
\NormalTok{summary[}\StringTok{"group"}\NormalTok{] }\OperatorTok{=}\NormalTok{ summary[}\StringTok{"group"}\NormalTok{].astype(}\BuiltInTok{str}\NormalTok{)}
\NormalTok{medians }\OperatorTok{=}\NormalTok{ summary.melt(}
\NormalTok{    id\_vars}\OperatorTok{=}\NormalTok{[}\StringTok{"group"}\NormalTok{, }\StringTok{"games"}\NormalTok{], value\_vars}\OperatorTok{=}\NormalTok{[}\StringTok{"actual"}\NormalTok{, }\StringTok{"predicted"}\NormalTok{],}
\NormalTok{    var\_name}\OperatorTok{=}\StringTok{"type"}\NormalTok{, value\_name}\OperatorTok{=}\StringTok{"reviews"}\NormalTok{,}
\NormalTok{)}
\NormalTok{medians[}\StringTok{"type"}\NormalTok{] }\OperatorTok{=}\NormalTok{ medians[}\StringTok{"type"}\NormalTok{].}\BuiltInTok{str}\NormalTok{.capitalize()}

\NormalTok{x\_group }\OperatorTok{=}\NormalTok{ alt.X(}\StringTok{"group"}\NormalTok{, sort}\OperatorTok{=}\VariableTok{None}\NormalTok{, title}\OperatorTok{=}\StringTok{"Actual reviews (grouped)"}\NormalTok{, axis}\OperatorTok{=}\NormalTok{alt.Axis(labelAngle}\OperatorTok{={-}}\DecValTok{45}\NormalTok{))}
\NormalTok{spread }\OperatorTok{=}\NormalTok{ alt.Chart(summary).mark\_rule(strokeWidth}\OperatorTok{=}\DecValTok{3}\NormalTok{, opacity}\OperatorTok{=}\FloatTok{0.4}\NormalTok{, color}\OperatorTok{=}\StringTok{"\#1f77b4"}\NormalTok{).encode(}
\NormalTok{    x}\OperatorTok{=}\NormalTok{x\_group,}
\NormalTok{    y}\OperatorTok{=}\NormalTok{alt.Y(}\StringTok{"predicted\_low"}\NormalTok{, scale}\OperatorTok{=}\NormalTok{scale, axis}\OperatorTok{=}\NormalTok{axis, title}\OperatorTok{=}\StringTok{"Number of reviews"}\NormalTok{),}
\NormalTok{    y2}\OperatorTok{=}\StringTok{"predicted\_high"}\NormalTok{,}
\NormalTok{)}
\NormalTok{dots }\OperatorTok{=}\NormalTok{ alt.Chart(medians).mark\_point(filled}\OperatorTok{=}\VariableTok{True}\NormalTok{, size}\OperatorTok{=}\DecValTok{60}\NormalTok{, opacity}\OperatorTok{=}\DecValTok{1}\NormalTok{).encode(}
\NormalTok{    x}\OperatorTok{=}\NormalTok{x\_group,}
\NormalTok{    y}\OperatorTok{=}\NormalTok{alt.Y(}\StringTok{"reviews"}\NormalTok{, scale}\OperatorTok{=}\NormalTok{scale, axis}\OperatorTok{=}\NormalTok{axis),}
\NormalTok{    color}\OperatorTok{=}\NormalTok{alt.Color(}\StringTok{"type"}\NormalTok{, title}\OperatorTok{=}\VariableTok{None}\NormalTok{,}
\NormalTok{                    scale}\OperatorTok{=}\NormalTok{alt.Scale(domain}\OperatorTok{=}\NormalTok{[}\StringTok{"Actual"}\NormalTok{, }\StringTok{"Predicted"}\NormalTok{], }\BuiltInTok{range}\OperatorTok{=}\NormalTok{[}\StringTok{"black"}\NormalTok{, }\StringTok{"\#1f77b4"}\NormalTok{])),}
\NormalTok{    tooltip}\OperatorTok{=}\NormalTok{[}
        \StringTok{"group"}\NormalTok{, }\StringTok{"type"}\NormalTok{,}
\NormalTok{        alt.Tooltip(}\StringTok{"reviews"}\NormalTok{, }\BuiltInTok{format}\OperatorTok{=}\StringTok{",.0f"}\NormalTok{, title}\OperatorTok{=}\StringTok{"Median reviews"}\NormalTok{),}
\NormalTok{        alt.Tooltip(}\StringTok{"games"}\NormalTok{, }\BuiltInTok{format}\OperatorTok{=}\StringTok{","}\NormalTok{, title}\OperatorTok{=}\StringTok{"Games in group"}\NormalTok{),}
\NormalTok{    ],}
\NormalTok{)}

\NormalTok{panel\_b }\OperatorTok{=}\NormalTok{ (spread }\OperatorTok{+}\NormalTok{ dots).properties(width}\OperatorTok{=}\DecValTok{300}\NormalTok{, height}\OperatorTok{=}\DecValTok{300}\NormalTok{, title}\OperatorTok{=}\StringTok{"B. Grouped by actual reviews"}\NormalTok{)}

\NormalTok{fig\_groups }\OperatorTok{=}\NormalTok{ (panel\_a }\OperatorTok{|}\NormalTok{ panel\_b).resolve\_scale(y}\OperatorTok{=}\StringTok{"shared"}\NormalTok{)}
\end{Highlighting}
\end{Shaded}

\begin{figure}

\centering{

\pandocbounded{\includegraphics[keepaspectratio]{videogames-regression_files/figure-pdf/fig-prediction-accuracy-output-1.png}}

}

\caption{\label{fig-prediction-accuracy}Predicted versus actual review
counts for test-set games. (A) Every test game, with the dashed line
marking perfect prediction and grey lines marking group boundaries. (B)
Median actual and predicted review counts within each group of actual
reviews; vertical bars show the middle 50\% of predictions.}

\end{figure}%

Figure~\ref{fig-prediction-accuracy} compares predicted and actual
review counts for the 9,372 test games in two ways. Panel A plots every
game. If the model were perfect, every point would sit on the dashed
line. The cloud is clearly flatter than the line, but it is hard to read
beyond that. With this many games the points overlap heavily. Review
counts are also whole numbers, so games with few reviews pile up in
vertical stripes at 0, 1, 2 and so on. From panel A alone it is hard to
tell where most games lie or what a typical prediction looks like at a
given level of popularity.

Panel B summarises the same games. It splits them into nine groups by
actual review count (the grey lines in panel A) and shows the median
actual and median predicted count in each group, along with the middle
50\% of predictions. Grouping has three advantages here:

\begin{itemize}
\tightlist
\item
  \textbf{No overplotting.} Each group is reduced to a few marks, so
  nothing is hidden behind other points.
\item
  \textbf{Equal weight for small groups.} The few hundred very popular
  games get as much space as the thousands of small ones, instead of
  being lost in a corner of the cloud.
\item
  \textbf{Robust summaries.} Medians are not pulled around by a handful
  of extreme games, and each group still holds between 285 and 1,945
  games, so the medians are stable.
\end{itemize}

Panel B shows that the model's predictions are pulled toward the middle.
Across the nine groups the median actual review count ranges from 0 to
8,261, but the median prediction only ranges from 6 to 226. As a result
the model over-predicts games with few or no reviews and under-predicts
popular ones. Games with more than 3,000 reviews are predicted to get
about 37 times fewer reviews than they actually received.

This pattern is expected from a model that explains only about a third
of the variance. Price, languages, genres and categories can separate a
polished game from a bare-bones one, but they cannot tell a hit from an
ordinary well-made game. The factors behind a breakout hit, such as
marketing, quality and word of mouth, are not in the data. The model is
therefore more useful for comparing games roughly than for predicting
the review count of any single game, especially a very popular one.

\section{Discussion}\label{discussion}

\subsection{Summary of Findings and Link to Research
Question}\label{summary-of-findings-and-link-to-research-question}

This study examined how well selected game characteristics predict the
total number of reviews received by Steam games released between 2018
and 2022. The results show that these features provide useful predictive
information, but explain only part of the differences in review counts,
accounting for approximately 34\% of the variation in log review counts
on the test set (R² = 0.340). Although the linear regression model
performed better than the mean baseline in cross-validation, its
predictions were less accurate for games with very few or very high
review counts.

\subsection{Real-world implications}\label{real-world-implications}

These findings suggest that game characteristics available on Steam
provide some useful information about review volume. Developers and
publishers may use these characteristics to compare games with different
features and understand general patterns in review activity. However,
the model should not be used to make precise predictions for individual
games. Review volume is also not a direct measure of game quality,
player satisfaction, or commercial success. Other factors, such as
marketing activities, community engagement, and external visibility, may
also influence review counts.

\subsection{Alignment with
Expectations}\label{alignment-with-expectations}

Some findings matched our expectations, including the strong right-skew
in review counts (mean = 1,158.54; median = 14) and higher median review
counts for games supporting more languages. The model also showed a
clear tendency to overpredict review counts for games with few or no
reviews and underpredict them for games with more than 3,000 reviews.

\subsection{Assumptions and
Limitations}\label{assumptions-and-limitations}

Our model struggles at both ends of the distribution of review counts. A
large portion of games in the training set have no reviews, and a linear
model using log1p(total\_reviews) treats these as low values rather than
as a separate group. Therefore, it does not explicitly model zero-review
games. As shown in Panel B of Figure~\ref{fig-prediction-accuracy},
predictions are pulled toward the middle: the model overpredicts review
counts for games with few or no reviews and underpredicts them for games
with the highest review counts. A model designed to account for excess
zero counts may help address this limitation. Linear regression assumes
that the relationships between the selected features and log review
counts can be represented by a linear combination. It may therefore miss
nonlinear relationships and interactions. The model also uses only five
groups of game characteristics and excludes potentially important
factors such as marketing activities, developer reputation, and
community size. In addition, games released in different years may have
had different amounts of time to accumulate reviews, affecting
comparisons. Finally, these results describe predictive relationships
and do not establish that any particular game characteristic causes an
increase in reviews.

\subsection{Potential Impact}\label{potential-impact}

These insights may help developers and publishers make more informed
decisions about game design and marketing by providing a general
understanding of the characteristics associated with review volume.
However, review counts should be considered alongside other indicators
of game performance.

\subsection{Future directions}\label{future-directions}

Future research could examine whether adding factors such as game
ratings, marketing activities, developer reputation, and community size
improves prediction performance. It could also investigate whether
relationships between game characteristics and review counts differ
across release years. Comparing linear regression with nonlinear,
zero-inflated, or two-part models may help address the prediction errors
observed in this study. Finally, applying similar methods to games from
other release periods or platforms could help assess whether these
findings generalize beyond the dataset.

\section{References}\label{references}

FronkonGames. (n.d.). Steam Games Dataset {[}Data set{]}. Hugging Face.
https://doi.org/10.57967/hf/0511

Guzsvinecz, T., \& Szűcs, J. (2026). Emotions in play: A length and
sentiment analysis of Steam reviews by player engagement mode. Heliyon,
12, e45258. https://doi.org/10.1016/j.heliyon.2026.e45258

Lin, D., Bezemer, C. P., Zou, Y., \& Hassan, A. E. (2019). An empirical
study of game reviews on the Steam platform. Empirical Software
Engineering, 24, 170--207. https://doi.org/10.1007/s10664-018-9627-4

Philp, M., \& Nepomuceno, M. V. (2024). How reviews influence product
usage post-purchase: An examination of video game playtime. Journal of
Business Research, 172, 114456.
https://doi.org/10.1016/j.jbusres.2023.114456

Valve. (n.d.). User reviews. Steamworks Documentation.
https://partner.steamgames.com/doc/store/reviews




\end{document}
